%% ============================================================
%%  00_abstract.tex – Abstract & Zusammenfassung
%% ============================================================
\chapter*{Abstract}
\addcontentsline{toc}{chapter}{Abstract}
% English summary (approx. 150–250 words).
This bachelor's thesis addresses the development and runtime analysis of an
abstraction layer that allows quantum circuits and observables to be described in a
single form and subsequently executed as objects of the quantum frameworks
\textit{PennyLane} and \textit{Qiskit}. The software component was integrated into
the \textit{Executor} Python library, which is currently under active development at
Fraunhofer IPA and is to be published under an open-source license in the near
future. Its previous internal representation relied on \textit{Qiskit} objects and
was therefore tied to a single framework. This dependency is replaced by a
self-contained data structure, so that the description of the circuit remains
unchanged and the execution environments of the respective other framework can be
addressed solely by selecting the backend. Since framework calls occur exclusively
within the backends, a program written once also remains executable across framework
version changes. In addition, the runtime overhead introduced by this layer compared
to direct use of the frameworks is evaluated.
\cleardoublepage

%% ---- Deutsche Zusammenfassung -----------------------------
\chapter*{Zusammenfassung}
\addcontentsline{toc}{chapter}{Zusammenfassung}
% Kurze Zusammenfassung auf Deutsch (ca. 150–250 Wörter).
Diese Bachelorarbeit befasst sich mit der Entwicklung und der Laufzeitanalyse einer
Abstraktionsschicht, über die sich Quantenschaltkreise und Observablen einheitlich
beschreiben und anschließend in den Objekten der Quantenframeworks
\textit{PennyLane} und \textit{Qiskit} ausführen lassen. Die Softwarekomponente
wurde in die Python-Softwarebibliothek \textit{Executor} integriert, die derzeit vom
Fraunhofer IPA aktiv weiterentwickelt wird und in naher Zukunft unter einer
Open-Source-Lizenz veröffentlicht werden soll. Bislang stützte sich deren interne
Darstellung auf \textit{Qiskit}-Objekte und war damit an ein einzelnes Framework
gebunden. Diese Bindung wird durch eine eigenständige Datenstruktur ersetzt, sodass
die Beschreibung des Schaltkreises unverändert bleibt und sich allein durch die
Auswahl des Backends die Ausführungsumgebungen des jeweils anderen Frameworks
ansprechen lassen. Da Framework-Aufrufe ausschließlich innerhalb der Backends
auftreten, bleibt ein einmal geschriebenes Programm zudem über Versionswechsel der
Frameworks hinweg lauffähig. Ergänzend wird der zeitliche Mehraufwand bewertet, den
diese Schicht gegenüber der direkten Nutzung der Frameworks verursacht.
\cleardoublepage